\paragraph{Theorem F.2.12 (Proved: an integral obstruction for a rational $(5,12)$ family).}
Let a rational rank-six cohomology representation have generators
\[
 A=\begin{pmatrix}1&0\\0&A_5\end{pmatrix},\qquad
 B=\begin{pmatrix}1&c\\0&B_5\end{pmatrix},\qquad
 c=(0,0,0,1,0),
\]
where
\[
 A_5=\begin{pmatrix}
 1&0&1&0&0\\1&0&0&0&0\\-1&0&0&1&1\\
 1&-1&0&-1&0\\0&1&0&0&0
 \end{pmatrix},\qquad
 B_5=\begin{pmatrix}
 0&1&0&0&0\\0&0&0&0&1\\1&-2&0&0&0\\
 0&1&0&-1&-1\\0&0&1&1&2
 \end{pmatrix}.
 \tag{F.33}
\]
For every invariant integral lattice with an index-one maximal
square-zero cusp, and every free affine filling of multiplicities
$m_1,m_2$ divisible respectively by five and twelve, a simply
connected total space has
\[
 \operatorname{Tor}_3 H_2
       \simeq\mathbb Z/3^{v_3(m_2)}.
 \tag{F.34}
\]
In particular none of these data is three-connected.
The standard lattice nevertheless has a primitive cusp,
free affine fillings of orders five and twelve, and a simply
connected total space with rational Betti vector
$(1,0,0,0,4,0,0,0,1)$.

\begin{proof}
All matrix identities and integral lattices used below are
printed and checked exactly in
\texttt{fourfolds\_512\_probe.py}.
The generators have orders five and twelve, characteristic
polynomials $(T-1)^2\Phi_5$ and $(T-1)^2\Phi_{12}$,
and $(AB-I)^2=0$ with three nonzero Smith factors equal to one.
Their common invariant line is the first coordinate line.
For homology matrices $A^t,B^t$ put
\[
 u_1=(1,0,0,0,0,0)^t,\qquad
 u_2=(1,-1,1,-1,0,-1)^t,
 \qquad \ell=(1,0,0,0,0,0).
\]
They are invariant vectors with $\ell(u_i)=1$.
The cusp vanishing lattice and its two finite translates
generate the saturated lattice $\ker\ell$.
Choose twists $v_1=-2u_1$, $v_2=5u_2$.
Their $\ell$-values are units modulo five and twelve, proving
every nonidentity affine power free.
After killing the vanishing lattice the group relations are
$5a=-2h$, $12b=5h$, $a+b=0$; their determinant has absolute
value $|-24+25|=1$.  Thus the total space is simply connected.
The exact rational exterior calculations give
\[
 h^0=(1,1,0,0,0,0,1),\qquad
 h^1=(0,0,0,3,0,0,0),\qquad
 h^2=(1,0,0,0,0,1,1)
\]
for the invariant sheaves, and the same rational Leray
argument as in F.2.5 gives the asserted Betti vector.

We first prove that changing the lattice cannot remove the
obstruction at the prime three.  Write $R=\mathbb Z_3$ and
$K=\mathbb Q_3$.  The standard quotient lattice $W_R$ defined
by $A_5,B_5$ has irreducible reduction modulo three.
Indeed $\Phi_5$ is irreducible over $\mathbb F_3$, since
three has order four modulo five.  Every proper nonzero
$A_5$-stable subspace therefore has dimension one or four.
On a one-dimensional subrepresentation or quotient its
$A_5$ character is one; the unipotent product forces the
$B_5$ character to be one too.  The stacked and horizontal
matrices for $A_5-I,B_5-I$ both have rank five modulo three,
as certified in the code, ruling out those two possibilities.

This irreducible reduction makes $W_R$ the unique stable
$R$-lattice up to homothety.  To see this, rescale another
lattice into $W_R$ but not into $3W_R$.  Its reduction is
a nonzero invariant subspace, hence all of $W_R/3W_R$.
Nakayama's lemma, or the elementary Smith form of the
inclusion, then makes it equal to $W_R$.

Normalize the quotient of any stable lattice in $V_K$ to
be $W_R$.  Its intersection with the invariant line is
$3^rRe$ for some integer $r$.
The representation of the free group on the two generators
has extension class
$\alpha\in H^1(\Gamma,W_R^*)$.
This cohomology group is torsion free: the coboundary map
from $W_R^*$ into two copies has full column rank modulo
three, since $W_R^*$ has no invariants modulo three.
The class $\alpha$ is primitive in this free $R$-module.
Otherwise the extension would split modulo three and $V$
would have a trivial quotient there.  The horizontal
matrix $[A-I\ B-I]$ instead has Smith factors
$(1,1,1,1,1,2)$, so it has full rank modulo three.

In a basis starting with $3^re$, the extension class is
$3^{-r}\alpha$.  Integrality is therefore possible exactly
when $r\leq0$.  Changes of splitting introduce no further
classes: the coboundary image is primitive, so an integral
coboundary over $K$ is already an integral coboundary over
$R$.  Thus these lattices are represented by multiplying
the extension rows by $3^{-r}$.
On the quotient cusp $N_W$ has primitive rank-two image.
The extra cusp map is the nonzero map
$\ker N_W/\operatorname{im}N_W\to Re$.
It is a unit for the displayed lattice, because its full
cusp has Smith factors $(1,1,1)$; for the lattice indexed
by $r$ it is multiplied by $3^{-r}$.
Hence the full cusp is primitive exactly when $r=0$.
We have proved uniqueness of the primitive-cusp lattice
over $R$, without relying on a bounded lattice search.

It remains to compute the obstruction on that lattice.
We first justify the arbitrary-twist comparison integrally.
For any global invariant lattice take the primitive invariant
homology covector $\ell$ and quotient the van Kampen group
by every fibre loop in $\ker\ell$.
The finite presentations become
$[a,h]=[b,h]=1$, $a^{m_1}=h^{r_1}$,
$b^{m_2}=h^{r_2}$, where $r_i=\ell(v_i)$.
The cusp meridian gives $ab=h^{-\kappa}$ for an integer
$\kappa$ determined by its translation clutching
($\kappa=0$ for the stated standard clutching).
The cusp angular lattice lies in $\ker\ell$, so it adds
no further relation to this quotient.  In particular this
is an abelian quotient with determinant
$m_1m_2\kappa+m_2r_1+m_1r_2$.
Simple connectivity makes that determinant $\pm1$.
Reduction modulo each $m_i$ proves both
$\gcd(m_1,m_2)=1$ and $\gcd(r_i,m_i)=1$.

The two finite fillings have free $R$-cohomology in the
degrees needed here.  In the invariant marking whose first
homology vector over $R$ has $\ell$-value one, the map
$m_i\ell$ from the reduced finite fibre to $\mathbb R/\mathbb Z$
has connected fibre $T^5$: the affine generator permutes
its $m_i$ preimage components by addition of $r_i$ modulo
$m_i$, transitively by the proved coprimality.
Its monodromy on vertical homology is $D_i^{s_i}$, where
$s_ir_i\equiv1\pmod {m_i}$, possibly followed by a translation.
Thus it is a mapping torus of a power prime to the linear
order.  The translation is isotopic to the identity and
does not change its Wang maps.  At the second point, over
$R$ we may take the invariant first vector to be $v_2/r_2$;
its denominator is prime to three and does not change the
$R$-homology lattice.  At the first point the covering
degree $m_1$ is prime to three, so averaging and the transfer
identify its $R$-cohomology with the invariant fibre
cohomology directly.
In degrees one and two the matrices
$D_i-I$ and $\bigwedge^2D_i-I$ have only unit Smith factors.
This proves the stated freeness by the Wang sequence.

The finite degree-two covering cokernel is read from
integral periods.  Vertical cycles have integral periods.
The remaining primitive two-cycle is the commuting torus
formed by a base suspension and a primitive fixed vertical
loop $\lambda_i$; its covering period is
$\omega(u_i,\lambda_i)/m_i$ for an invariant integral
two-form $\omega$.
More explicitly the homological Wang sequence is
\[
 0\longrightarrow\operatorname{coker}(D_i^{s_i}-I\mid H_2(T^5;R))
 \longrightarrow H_2(B_i;R)
 \longrightarrow\ker(D_i^{s_i}-I\mid H_1(T^5;R))
 \longrightarrow0.
\]
The last group has rank one, generated by $\lambda_i$.
The commuting suspension torus maps to that generator;
at the second point the invariant vector $v_2/r_2$ may be used after taking
$R$-linear combinations of integral cycles.
Therefore vertical cycles and this one torus generate
$H_2(B_i;R)$, proving sufficiency of the period test.
For integral bases of the invariant two-forms, the numerator
rows at the two points are
\[
 (0,0,5),\qquad (3,1,1).
 \tag{F.35}
\]
The vertical Wang quotient has no torsion, so those cycles
generate the full free second homology and the period test
is sufficient as well as necessary.

Since $3\mid m_2$, the already-proved coprimality gives
$3\nmid m_1$ and makes $r_2$ a unit in $R$.
Changing that unit twist or adding a fixed vertical vector
does not change the three-primary cokernel in (F.35): the
former multiplies its generating period by a unit, and the
latter has zero wedge with the fixed vertical loop.
The first finite point has zero covering cokernel over $R$;
the second has $R/m_2R$.

The global invariant degree-two group is zero, and the
degree-two global coboundary image is primitive, with
vanishing rational parabolic first cohomology.  Therefore
$H^1(\mathbb P^1,R^2g_*R)=R/m_2R$.
The global degree-one coinvariants vanish over $R$, and
there are no degree-two or degree-three point-supported
kernels: the finite cap groups just proved free, and the
Kuhn cusp has its sole kernel in degree four.
The degree-three global invariant group also vanishes.
Consequently the only nonzero total-degree-three Leray
term over $R$ is $E_2^{1,2}=R/m_2R$; it has no possible
differential.  Thus $H^3(g^{-1}(\mathbb P^1);R)=R/m_2R$.
The universal coefficient theorem gives precisely (F.34).
\end{proof}
